\paragraph{F.46 (Proved). The point function of every cyclotomic test morphism.}
Let $k=\mathbb Q(\mu_{2^\infty})$, fix $K=\overline{\mathbb Q}$, and let
\[
 \phi:(\mathbb G_{m,k})_{\mathrm{et}}
       \longrightarrow(\mathbb G_{m,k})_{\mathrm{et}}
\]
be a pinned geometric morphism over $B G_k$, with the relative
identification included in the data. Its induced map on geometric
points is
\[
 f:K^*\longrightarrow K^*,\qquad f(a)=c a^m
 \quad(c\in k^*,\ m\in\mathbb Z).
 \tag{F46}
\]
Thus the point-function obstruction in the cyclotomic test over
$\mathbb Q(\mu_{2^\infty})$ vanishes. This assertion by itself does not identify the full geometric
morphism: the hypothesis of Carlson--Haine--Wolf,
Proposition~4.14, concerns every quasicompact separated etale target
object and its pullback, not only the terminal object.

\emph{Problem locus.} Montpellier cote149, \S2, fol.25v, scan52,
and the supplied prompt's Part~F intermediate-field reconstruction question.
This develops the arithmetic reconstruction question there.
It is a new consequence of the proved elementary arguments
below and the explicitly cited arithmetic input; it is not attributed
to a sentence of the manuscript.

\emph{Proof: arithmetic facts.}
For a field $L$ put $D(L^*)=\bigcap_{n\geq1}(L^*)^n$.
For every finite extension $L/k$ one has
\[
 D(L^*)=\mu_{2^\infty}.                                      \tag{F46.1}
\]
Here is the precise arithmetic input and its application.
Amoroso--Zannier, Theorem~1.1, gives a positive lower bound for the
height of a non-torsion algebraic number in terms of its degree over
an abelian extension of a fixed number field
\cite{AmorosoZannier2000}. Apply it with the abelian extension
$\mathbb Q^{\mathrm{ab}}/\mathbb Q$. The degrees of elements of
$L\mathbb Q^{\mathrm{ab}}$ over $\mathbb Q^{\mathrm{ab}}$ are bounded
by $[L:k]$. Consequently the heights of non-torsion elements of $L^*$
have a uniform positive lower bound. If $a\in D(L^*)$ were not a root
of unity, choosing $b_n\in L$ with $b_n^n=a$ would give
$h(b_n)=h(a)/n>0$, a contradiction. Thus $a$ is torsion.

The field $k$ is unramified at each odd prime $\ell$: every finite
subextension is contained in a $2$-power cyclotomic extension, whose
defining roots of unity are etale over $\mathbb Z_\ell$. The value
group of any extension of the $\ell$-adic valuation to $L$ therefore
has index at most $[L:k]$ over $\mathbb Z$. In particular, the orders
of the $\ell$-primary roots of unity in $L$ are bounded, since
\[
 v_\ell(\zeta_{\ell^r}-1)
       =\frac{1}{\ell^{r-1}(\ell-1)}.
\]
If a torsion element of $D(L^*)$ had nontrivial $\ell$-primary part,
its $\ell^n$-th roots in $L$ would contradict this bound. Only its
$2$-primary part remains. Conversely $\mu_{2^\infty}$ is divisible
by every positive integer and is contained in $k$, proving (F46.1).

We also use that $K/k$ has infinite degree, and that every finite
extension of $k$ admits a discrete valuation. Choose a valuation over
$3$ with $v(3)=1$. Its value group on $k$ is $\mathbb Z$; on a finite
extension it has finite index over $\mathbb Z$ and can be renormalized
to $\mathbb Z$. Finally $T^n-3$ is Eisenstein for this valuation, so
$[K:k]$ is unbounded. Equivalently, a root has valuation $1/n$, which
forces the degree of its extension of $k$ to be at least $n$.

\emph{Proof: a rational-source Kummer lemma.}
Let $L/k$ be finite, and let
\[
 U=\mathbb P^1_L\setminus\{a_1,\ldots,a_r,\infty\},
 \qquad a_i\in L\text{ distinct}.
\]
Suppose a compatible system of Kummer classes
$\alpha\in\varprojlim_n H^1(U,\mu_n)$ has the following property:
at every $x\in U(L')$, for every finite $L'/L$, its evaluation is the
Kummer class of a specified element $F(x)\in(L')^*$.
Then there is a regular unit $g\in\mathcal O(U)^*$ such that
\[
 \alpha=[g],\qquad F(x)/g(x)\in D((L')^*)
       \quad(x\in U(L')).                                  \tag{F46.2}
\]
Indeed, $\operatorname{Pic}(U)=0$, and unique factorization gives
\[
 \alpha=\widehat c\prod_{i=1}^r(t-a_i)^{\lambda_i},\qquad
 \widehat c\in\widehat{L^*},\quad\lambda_i\in\widehat{\mathbb Z}.
\]
Choose a discrete valuation $v:L^*\to\mathbb Z$ and an element
$\pi$ of value one. For fixed $j$, put $x_N=a_j+\pi^N$.
For all sufficiently large integers $N$, these are points of $U(L)$,
$v(x_N-a_j)=N$, and
$v(x_N-a_i)=v(a_j-a_i)$ for $i\ne j$.
Evaluating each Kummer quotient and subtracting the equations for
$N+1$ and $N$ yields
\[
 \lambda_j=v(F(x_{N+1}))-v(F(x_N))
       \quad\text{in }\widehat{\mathbb Z}.
\]
Its right side is an integer, so write $\lambda_j=m_j\in\mathbb Z$.
Choose $x_0\in U(L)$ and set
\[
 c=F(x_0)\Big/\prod_i(x_0-a_i)^{m_i}.
\]
Evaluation at $x_0$ identifies $\widehat c$ with the completion of
this actual $c\in L^*$. Thus $g=c\prod_i(t-a_i)^{m_i}$ proves
(F46.2), including its assertion after every finite extension.
No finite-presentation descent for the topos morphism is used here.

\emph{Proof: apply the lemma to every translate.}
A relative pinned morphism sends $L$-rational points to
$L$-rational points, for every finite $L/k$. One way to check this
is to compose the point's section $G_L\to\pi_1$ with the given map.
If its image is a closed point with residue field $E$, the resulting
map into $G_E$ lies over the inclusion $G_L\subset G_k$; hence
$G_L\subset G_E$ and the specified embedding of $E$ lies in $L$.
The induced function $f$ is therefore $G_k$-equivariant and has the
stated finite-extension compatibility. In particular $f(1)\in k^*$;
divide by this constant and assume from now on that $f(1)=1$.

If the generic point maps to a closed point, continuity and pinning
make $f$ constant, necessarily equal to $1$, and (F46) follows with
$m=0$. Otherwise its image is generic. For each $b\in K$, after a
finite extension containing $b$, pull back the open
$\mathbb G_m\setminus\{b\}$ of the target. This is a nonempty
open $U_b$ of the source, with finitely many missing points. Its
geometric points are exactly the $x$ with $f(x)\ne b$.
After a further finite extension $L/k$ all the missing points are
rational. Pullback of the compatible Kummer classes of $t-b$ to
$U_b$ satisfies the preceding lemma, with $F(x)=f(x)-b$.
Thus there is $g_b\in L(t)^*$, a unit on $U_b$, such that
\[
 f(x)-b=d_b(x)g_b(x),\qquad d_b(x)\in\mu_{2^\infty}\subset k^*,
       \quad x\in U_b(K).                                  \tag{F46.3}
\]
Each $x$ is contained in a finite extension of $L$, so (F46.1)
applies. For $b=0$ the domain is all of $\mathbb G_m$. The lemma
and the normalization at $1$ give
\[
 f(x)=d_0(x)x^m,\qquad m\in\mathbb Z,\quad d_0(x)\in k^*.
                                                               \tag{F46.4}
\]
The stronger assertion $d_0(x)\in\mu_{2^\infty}$ will not be needed
until the end.

\emph{Proof: the case $m=0$.}
In this case all values of $f$ lie in $k$. Choose $b\in K\setminus k$.
Then $U_b=\mathbb G_m$, and (F46.3) gives
$f(x)-b=d_b(x)c_bx^n$ for an integer $n$ and $c_b\in K^*$.
If $n\ne0$, the values of $c_bx^n$ fill $K^*$; the equality would
place all these values in the $k$-vector space $k+kb$, of dimension
at most two. This contradicts $[K:k]=\infty$.
Therefore $n=0$. If $f$ took two distinct values, subtracting their
two equations would give $c_b\in k$ and then $b\in k$, again a
contradiction. So $f$ is constant, and $f(1)=1$ gives $f=1$.
This treats the zero exponent without applying a nonconstant
quasi-rationality criterion to it.

\emph{Proof: a finite-exception quasi-rationality lemma.}
We give the argument needed to use (F46.3) safely at its missing
points. It is the finite-exception form of the mechanism in
Tsujimura, Lemmas~3.3--3.5 and Proposition~3.6
\cite{Tsujimura2022}.
Suppose $K$ is algebraically closed of characteristic zero, algebraic over a subfield $M$,
and $[K:M]=\infty$. Let $C/K$ be a smooth proper connected curve,
let $q$ be a function on a cofinite subset of $C(K)$, and suppose
that for each $b\in K$ one has
\[
 q(x)-b=e_b(x)r_b(x),\qquad e_b(x)\in M^*,\quad r_b\in K(C)^*,
                                                               \tag{F46.5}
\]
outside a finite set depending on $b$. Then $q$ agrees with one
rational function on a cofinite subset. If $r_0$ is constant, this
rational function is constant as well.

First suppose $r_0$ is nonconstant and apply (F46.5) to $q$ and
$1-q$. Write their rational factors as $r,s$. Choose a finite
extension $N/M$ in $K$ over which $C,r,s$ and a finite nonconstant
map $C\to\mathbb P^1$ are defined. For fixed $A,B\in M^*$ the
degree of the pole divisor of $Ar+Bs-1$, if it is nonzero, is
bounded by one integer $d$, independent of $A,B$. Thus a zero
of this nonzero rational function has residue degree at most $d$
over $N$. Choose $x$ outside the exceptional sets whose residue
field over $N$ has degree greater than $d$. Such points exist:
choose a value of the fixed finite map of degree greater than
$d$ over $N$, and any point above it. The values, and hence the
points, can avoid any given finite set because $K/N$ has infinite
degree. Substituting $A=e_0(x)$ and the corresponding coefficient
$B$ of $1-q$ shows that $Ar+Bs-1$ vanishes at $x$.
The degree bound forces the rational-function identity $Ar+Bs=1$.
Rescale $r,s$ by these constants so that $r+s=1$.
The pointwise equation then gives
\[
 (e_0(x)-e_1(x))r(x)=1-e_1(x).
\]
Consequently $q(x)=r(x)$ wherever $r(x)\notin M$, outside the
finite exceptional set.

Choose $a\in K\setminus M$ and repeat this argument for
$q+a$ and $1-q-a$. Its rational factor $r_a$ is nonconstant:
if it were a constant $c$, the equation
$q+a=e_a c$, together with $q=e_0r_0$, would place all cofinite
values of the nonconstant rational function $r_0$ in the
finite-dimensional $M$-vector space $Mc+Ma$. A nonconstant
rational function has cofinite image, so this is impossible.
After rescaling we obtain $q(x)=r_a(x)-a$ wherever
$r_a(x)\notin M$.
The two sets $r(x)\notin M$ and $r_a(x)\notin M$ cover the
cofinite domain under consideration: if both rational values were
in $M$, (F46.5) would put both $q(x)$ and $q(x)+a$ in $M$.
Their intersection is infinite. To check this last point without
a cardinality assumption, enlarge $N$ by their finitely many
coefficients. Points of sufficiently large degree over $N$ have
neither rational value in $M$, since a fibre of either fixed
nonconstant rational function has bounded degree. Infinitely many
such points exist. On this intersection $r=r_a-a$; equality of
rational functions follows. Their common value therefore agrees
with $q$ outside a finite set.

If $r_0=c$ is a nonzero constant, divide $q$ by $c$; the hypothesis
for all translates is preserved, and $q(x)\in M$ outside a finite
set. Choose $b\in K\setminus M$. Its rational factor $r_b$ cannot
be nonconstant, because its cofinite image would be contained in
$M+Mb$ by (F46.5). Hence $r_b$ is constant. If $q$ took two distinct
values on the common cofinite domain, subtracting the equations
$q(x)-b=e_b(x)r_b$ would put $r_b$ and then $b$ in $M$.
Thus $q$ is constant there. This proves the lemma in both cases,
for an arbitrary smooth proper curve, as needed also in F.47.

\emph{Proof: remove the exceptional values and finish.}
For $m\ne0$, apply the lemma with $M=k$, $q=f$, and $r_0=t^m$.
Let $R\in K(t)$ be the resulting nonconstant rational function,
equal to $f$ outside a finite set. For each $b$, equation (F46.3)
implies that
\[
 (R-b)/g_b
\]
takes values in $k^*$ on a cofinite set. A nonconstant rational
function has cofinite image in $K$, whereas $k$ is a proper
subfield; hence this quotient is a constant $c_b\in k^*$.
Thus $R-b=c_b g_b$ as rational functions.

Fix any $x\in K^*$. Taking $b\ne f(x)$ ensures $x\in U_b$,
where $g_b$ is a regular unit. The rational identity just proved
first shows that $R$ has no pole at $x$. For every
$b\notin\{f(x),R(x)\}$ it then gives
\[
 \frac{f(x)-b}{R(x)-b}\in k^*.
\]
If $f(x)\ne R(x)$, the displayed fractional linear expression,
as $b$ varies, takes every value in $K\setminus\{0,1\}$.
This contradicts $k\ne K$. Therefore $f(x)=R(x)$ for every
$x\in K^*$, including the initially exceptional points.
In particular $R$ is a regular unit on $\mathbb G_{m,K}$, so
$R=c t^j$. Alternatively, (F46.4) shows directly that
$R/t^m$ is a rational function whose cofinite values lie in $k$;
it must be constant. Normalization at $1$ makes that constant one.
Undoing the initial normalization gives exactly (F46).
\hfill$\square$
